\section{Control FOC}

El FOC constituye el controlador principal del PMSM. Su función es transformar las corrientes trifásicas al marco síncrono $dq$, regular los componentes de corriente mediante lazos PI y generar las referencias de tensión que permiten actuar sobre la planta.

\subsection{Flujo algorítmico}
\begin{equation}
(i_a,i_b,i_c)\rightarrow\mathrm{Clarke}\rightarrow(i_\alpha,i_\beta)\rightarrow\mathrm{Park}\rightarrow(i_d,i_q)\rightarrow\mathrm{PI}\rightarrow(v_d,v_q)\rightarrow\mathrm{Park}^{-1}\rightarrow\mathrm{M3/M2}\rightarrow\mathrm{PMSM}.
\end{equation}

\subsection{Lazo de velocidad}
\begin{equation}
e_\omega=\omega^*-\omega_m,
\end{equation}
\begin{equation}
i_{q,raw}^*=K_{p\omega}e_\omega+K_{i\omega}\int e_\omega\,dt,
\end{equation}
\begin{equation}
i_q^*=\operatorname{sat}(i_{q,raw}^*,i_{q,min},i_{q,max}).
\end{equation}

La actualización del lazo de velocidad se realiza cada cuatro pasos del lazo de corriente: $T_{s,current}=50\,\mu s$ y $T_{s,speed}=200\,\mu s$.

\subsection{Lazo de corriente}
Las referencias son $i_d^*=0$ e $i_q^*$ proveniente del lazo de velocidad:
\begin{align}
e_d &= i_d^*-i_d,\\
e_q &= i_q^*-i_q.
\end{align}

Los controles antes de saturación son:
\begin{align}
v_{d,raw} &= K_{pd}e_d+K_{id}\int e_d\,dt+v_{d,dec},\\
v_{q,raw} &= K_{pq}e_q+K_{iq}\int e_q\,dt+v_{q,dec}.
\end{align}

\subsection{Desacoplamiento}
\begin{align}
v_{d,dec}&=-\omega_eL_qi_q,\\
v_{q,dec}&=\omega_e(L_di_d+\lambda_m).
\end{align}

\subsection{Saturación y anti-windup}
\begin{equation}
V_{max}=0.95\frac{V_{DC}}{\sqrt{3}}.
\end{equation}
La implementación reconstruida revierte el incremento integral correspondiente cuando se detecta saturación.

\subsection{Parametrización recuperada del FOC/M6}
\begin{table}[H]
\centering
\caption{Parámetros del controlador FOC reconstruido.}
\begin{tabular}{lll}
\toprule
Parámetro & Valor & Capa \\
\midrule
$L_d,L_q$ & \SI{0.0012}{H} & FOC \\
$\lambda_m$ & \SI{0.1100}{Wb} & Referencia histórica M6 \\
$p$ & 4 & FOC \\
$K_{p\omega}$ & 5.2804 & FOC \\
$K_{i\omega}$ & 1.2693 & FOC \\
$K_{pd},K_{pq}$ & 15.0796 & FOC \\
$K_{id},K_{iq}$ & 628.3185 & FOC \\
$i_{q,max/min}$ & $\pm112.43$ A & FOC \\
\bottomrule
\end{tabular}
\end{table}

La diferencia entre $\lambda_m=0.107245$ Wb en la planta y $\lambda_m=0.1100$ Wb en el FOC recuperado se mantiene explícita hasta cerrar su procedencia.
